\section{Gauge and circuit holonomy}

An equivariant change of fiber coordinates is specified by
$g=(g_0,\ldots,g_{n-1})\in R^n$ and acts by
\[
G_g(i,x)=(i,xg_i).
\]
We use the conjugation convention
\[
T_{v'}=G_gT_vG_g^{-1}.
\]
It produces the transformed voltages
\[
v_i'=g_i^{-1}v_i g_{i+1},
\]
with indices read modulo $n$.

\begin{definition}
The circuit holonomy is the ordered product
\[
h(v)=v_0v_1\cdots v_{n-1}\in R.
\]
\end{definition}

The transformed product telescopes:
\[
h(v')=g_0^{-1}h(v)g_0.
\]
Thus gauge does not preserve a chosen holonomy element in a nonabelian group;
it preserves its conjugacy class. In an abelian group, conjugacy is trivial
and the element itself is gauge invariant.

\begin{proposition}[One-edge normalization]
Every voltage assignment is gauge equivalent to
\[
(e,e,\ldots,e,h),
\]
where $h$ is its circuit holonomy.
\end{proposition}

\begin{proof}
Set $g_0=e$ and recursively choose
\[
g_{i+1}=v_i^{-1}g_i
\]
for $0\leq i<n-1$. Then $g_i^{-1}v_i g_{i+1}=e$ on the first $n-1$
edges. On the final edge, the transformed voltage is
\[
g_{n-1}^{-1}v_{n-1}g_0
=v_0v_1\cdots v_{n-1}=h.
\]
\end{proof}

